%! AP Statistics — Unit 1, Lesson 1.9 — Slides (Beamer)
\documentclass[aspectratio=169, 11pt]{beamer}
\usepackage{apstats-beamer}

%=============================================================================
\begin{document}

% ─────────────────────────────────────────────────────────────────────────────
% SLIDE 1 — LESSON TITLE
% ─────────────────────────────────────────────────────────────────────────────
{
\setbeamercolor{background canvas}{bg=navy}
\begin{frame}[plain]
  \begin{tikzpicture}[remember picture, overlay]
    \fill[navylight] (current page.north west)
      rectangle ([yshift=-1.35cm]current page.north east);
  \end{tikzpicture}
  \vspace{2.8cm}
  \hspace{0.55cm}%
  \begin{minipage}{0.9\textwidth}
    {\color{goldacc}\bfseries\small LESSON 1.9}\\[0.5cm]
    {\color{white}\bfseries\LARGE Comparing Distributions of\\[0.25cm]
      a Quantitative Variable}\\[1.4cm]
    {\color{skymid}\small AP Statistics~$\cdot$~Unit 1: Exploring One-Variable Data}
  \end{minipage}
\end{frame}
}

% ─────────────────────────────────────────────────────────────────────────────
% SLIDE 2 — WARM-UP
% ─────────────────────────────────────────────────────────────────────────────
{
\setbeamercolor{background canvas}{bg=sky}
\begin{frame}[t, plain]
  \navyheader{Warm-Up}
  \vspace{1em}%
  \hspace{0.55cm}%
  \begin{minipage}{0.9\textwidth}
    \small
    \textit{Use the parallel boxplots on your warm-up sheet showing ACT scores
    for No Prep vs.\ Test Prep students.}

    \vspace{0.5em}
    \begin{enumerate}
      \setlength{\itemsep}{0.6em}\setlength{\topsep}{0em}
      \item What is the median for each group? Which group scored higher?
      \item Find the IQR for each group. Which is more spread out?
      \item Describe the overlap between the two distributions.
      \item Write one AP comparative sentence about center.
            Name both groups; use a comparison word; include context.
    \end{enumerate}
    \vspace{0.4cm}
    {\color{navy}\itshape\small
      Today we learn to go beyond individual descriptions ---
      we compare two groups on the same display.}
  \end{minipage}
\end{frame}
}

% ─────────────────────────────────────────────────────────────────────────────
% SLIDE 3 — PRIMARY OBJECTIVE
% ─────────────────────────────────────────────────────────────────────────────
{
\setbeamercolor{background canvas}{bg=sky}
\begin{frame}[t, plain]
  \begin{tikzpicture}[remember picture, overlay]
    \fill[navy] (current page.north west)
      rectangle ([yshift=-0.25cm]current page.north east);
  \end{tikzpicture}
  \vspace{0.45cm}\par
  \hspace{0.55cm}%
  \begin{minipage}{0.9\textwidth}

    \sectionlabel{Primary Objective}

    \normalsize
    Students will \textbf{compare distributions} of a quantitative variable
    using parallel boxplots and back-to-back stemplots, and write AP-quality
    comparative statements that address center, spread, shape, and outliers.

    \vspace{0.5cm}
    \textbf{Learning Targets:}
    \begin{enumerate}\small
      \setlength{\itemsep}{0.4em}
      \item Read parallel boxplots: identify five-number summaries for each group.
      \item Compare centers and spreads using explicit comparative language.
      \item Read a back-to-back stemplot; identify centers and shapes.
      \item Write a complete SOCS comparison that names both groups with context.
    \end{enumerate}
  \end{minipage}
\end{frame}
}

% ─────────────────────────────────────────────────────────────────────────────
% SLIDE 4 — HOOK: WHY SAME SCALE MATTERS
% ─────────────────────────────────────────────────────────────────────────────
{
\setbeamercolor{background canvas}{bg=goldbg}
\begin{frame}[t, plain]
  \navyheader{Hook: Why Does Same Scale Matter?}
  \vspace{0.8em}%
  \hspace{0.55cm}%
  \begin{minipage}{0.9\textwidth}
    \small
    ACT scores for No Prep vs.\ Test Prep students, placed on \textbf{one axis}:

    \vspace{0.5em}
    \begin{center}
    \begin{tikzpicture}[scale=1.05]
      \draw[thick] (0.0, 0) -- (6.80, 0);
      \foreach \v/\x in {14/0.00, 16/0.60, 18/1.20, 20/1.80, 22/2.40,
                          24/3.00, 26/3.60, 28/4.20, 30/4.80, 32/5.40, 34/6.00, 36/6.60} {
        \draw (\x, -0.12) -- (\x, 0.12);
        \node[below, font=\small] at (\x, -0.12) {\v};
      }
      \node[below, font=\small\itshape] at (3.30, -0.56) {ACT composite score};
      \node[left, font=\small] at (0.0, 1.10) {No Prep};
      \draw[thick, navy] (0.30, 1.10) -- (1.50, 1.10);
      \draw[thick, navy] (1.50, 0.85) rectangle (3.30, 1.35);
      \draw[thick, navy] (2.40, 0.85) -- (2.40, 1.35);
      \draw[thick, navy] (3.30, 1.10) -- (4.50, 1.10);
      \draw[thick, navy] (0.30, 0.95) -- (0.30, 1.25);
      \draw[thick, navy] (4.50, 0.95) -- (4.50, 1.25);
      \node[left, font=\small] at (0.0, 0.35) {Test Prep};
      \draw[thick, navy] (1.80, 0.35) -- (3.00, 0.35);
      \draw[thick, navy] (3.00, 0.10) rectangle (4.80, 0.60);
      \draw[thick, navy] (3.90, 0.10) -- (3.90, 0.60);
      \draw[thick, navy] (4.80, 0.35) -- (6.30, 0.35);
      \draw[thick, navy] (1.80, 0.20) -- (1.80, 0.50);
      \draw[thick, navy] (6.30, 0.20) -- (6.30, 0.50);
    \end{tikzpicture}
    \end{center}

    \vspace{0.3em}
    \begin{itemize}
      \setlength{\itemsep}{0.4em}
      \item At a glance: which group scored \textbf{higher}? Which is more \textbf{consistent}?
      \item Is there a meaningful \textbf{overlap}? What does that tell us?
      \item Could you answer these questions if the boxplots were on \textit{separate} scales?
    \end{itemize}
  \end{minipage}
\end{frame}
}

% ─────────────────────────────────────────────────────────────────────────────
% SLIDE 5 — READING PARALLEL BOXPLOTS
% ─────────────────────────────────────────────────────────────────────────────
{
\setbeamercolor{background canvas}{bg=white}
\begin{frame}[t, plain]
  \begin{tikzpicture}[remember picture, overlay]
    \fill[navy] (current page.north west)
      rectangle ([yshift=-0.25cm]current page.north east);
  \end{tikzpicture}
  \vspace{0.45cm}\par
  \hspace{0.55cm}%
  \begin{minipage}{0.9\textwidth}
    \sectionlabel{Reading Parallel Boxplots}
    \small
    \begin{columns}[T]
    \begin{column}{0.45\textwidth}
      \textbf{For each group, read:}
      \begin{itemize}\setlength{\itemsep}{0.35em}
        \item Median (line inside box)
        \item $Q_1$, $Q_3$, IQR (box width)
        \item Min, max, range
        \item Outlier dots (if any)
        \item Approximate skewness
      \end{itemize}

      \vspace{0.5em}
      \textbf{From the ACT display:}\\[0.3em]
      \begin{tabular}{lcc}
        & No Prep & Test Prep \\
        Median & 22 & 27 \\
        IQR    & 6  & 6  \\
        Range  & 14 & 15 \\
      \end{tabular}
    \end{column}
    \begin{column}{0.50\textwidth}
      \textbf{Across groups, ask:}
      \begin{itemize}\setlength{\itemsep}{0.4em}
        \item Which median is higher? By how much?
        \item Which IQR is larger? What does that mean?
        \item How much do the IQR boxes overlap?\\
              Large overlap $\Rightarrow$ groups are similar\\
              Small overlap $\Rightarrow$ real difference
        \item Any outliers in one group but not the other?
      \end{itemize}

      \vspace{0.4em}
      {\color{redacc}\textbf{Requirement:} Both boxplots on
      the \textit{same} number line. Separate scales
      make comparison impossible.}
    \end{column}
    \end{columns}
  \end{minipage}
\end{frame}
}

% ─────────────────────────────────────────────────────────────────────────────
% SLIDE 6 — AP COMPARATIVE LANGUAGE
% ─────────────────────────────────────────────────────────────────────────────
{
\setbeamercolor{background canvas}{bg=sky}
\begin{frame}[t, plain]
  \navyheader{AP Comparative Language}
  \vspace{0.8em}%
  \hspace{0.55cm}%
  \begin{minipage}{0.9\textwidth}
    \small
    \begin{columns}[T]
    \begin{column}{0.45\textwidth}
      \textbf{AP scoring rule:} A comparative statement earns credit only if it:
      \begin{enumerate}\setlength{\itemsep}{0.4em}
        \item[a.] Uses a \textbf{comparison word}
                  (higher, lower, more, less, similar, \ldots)
        \item[b.] \textbf{Names both groups} with context
        \item[c.] Gives \textbf{specific values} for both
      \end{enumerate}

      \vspace{0.6em}
      {\color{redacc}\textbf{No credit:}}\\
      ``The median for Test Prep was 27.''\\[0.3em]
      \textit{Problem: only one group; no comparison.}

      \vspace{0.4em}
      {\color{navy}\textbf{Full credit:}}\\
      ``The median ACT score for the Test Prep group
      (27) was \textbf{5 points higher than} the median
      for the No Prep group (22).''
    \end{column}
    \begin{column}{0.50\textwidth}
      \textbf{The comparison sentence frame:}

      \vspace{0.4em}
      \setbeamercolor{block body}{bg=goldbg, fg=navy}
      \setbeamercolor{block title}{bg=goldacc, fg=white}
      \begin{block}{}
      \small
      ``The \textbf{median} [variable] for [Group A]
      (\textit{value}) was \textbf{[higher/lower] than}
      the median for [Group B] (\textit{value}).''
      \end{block}

      \vspace{0.5em}
      \textbf{Full comparison = SOCS:}
      \begin{itemize}\setlength{\itemsep}{0.3em}
        \item Shape: describe both
        \item Outliers: note and verify
        \item Center: compare medians
        \item Spread: compare IQRs (both required)
      \end{itemize}
    \end{column}
    \end{columns}
  \end{minipage}
\end{frame}
}

% ─────────────────────────────────────────────────────────────────────────────
% SLIDE 7 — BACK-TO-BACK STEMPLOTS
% ─────────────────────────────────────────────────────────────────────────────
{
\setbeamercolor{background canvas}{bg=white}
\begin{frame}[t, plain]
  \begin{tikzpicture}[remember picture, overlay]
    \fill[navy] (current page.north west)
      rectangle ([yshift=-0.25cm]current page.north east);
  \end{tikzpicture}
  \vspace{0.45cm}\par
  \hspace{0.55cm}%
  \begin{minipage}{0.9\textwidth}
    \sectionlabel{Back-to-Back Stemplots}
    \small
    \begin{columns}[T]
    \begin{column}{0.45\textwidth}
      \textbf{When to use:}
      \begin{itemize}\setlength{\itemsep}{0.3em}
        \item Small datasets ($n < 30$)
        \item Want to see individual values and shape
        \item Complements parallel boxplots
      \end{itemize}

      \vspace{0.4em}
      \textbf{Structure:}
      \begin{itemize}\setlength{\itemsep}{0.3em}
        \item Shared stem column in the middle
        \item Left group: leaves in \textit{reverse} order
              (read outward, smallest to largest)
        \item Right group: leaves in normal order
      \end{itemize}

      \vspace{0.4em}
      \textbf{Reading:}
      \begin{itemize}\setlength{\itemsep}{0.3em}
        \item Median: count to middle leaf in each half
        \item Shape: where do leaves cluster?
        \item Compare tails and centers visually
      \end{itemize}
    \end{column}
    \begin{column}{0.50\textwidth}
      \textbf{ACT example} (key: $2\mid4 = 24$):

      \vspace{0.5em}
      \begin{center}
      \renewcommand{\arraystretch}{1.3}
      \begin{tabular}{r c l}
      \textit{No Prep (rev.)} & Stem & \textit{Test Prep} \\
      \hline
      8\;7 & 1 & \\
      8\;7\;5\;3\;2\;0 & 2 & 2\;4\;5\;6\;7\;8\;9 \\
       & 3 & 0\;3 \\
      \hline
      \end{tabular}
      \end{center}

      \vspace{0.5em}
      \begin{itemize}\setlength{\itemsep}{0.3em}
        \item No Prep median $= 23$ (4th of 8 values)
        \item Test Prep median $= 27$ (5th of 9 values)
        \item Test Prep leaves extend further right:\\
              higher center and higher maximum (33)
      \end{itemize}
    \end{column}
    \end{columns}
  \end{minipage}
\end{frame}
}

% ─────────────────────────────────────────────────────────────────────────────
% SLIDE 8 — GUIDED NOTES CHECK
% ─────────────────────────────────────────────────────────────────────────────
{
\setbeamercolor{background canvas}{bg=goldbg}
\begin{frame}[t, plain]
  \navyheader{Quick Check --- Guided Notes}
  \vspace{0.8em}%
  \hspace{0.55cm}%
  \begin{minipage}{0.9\textwidth}
    \small
    Using the ACT parallel boxplots (No Prep vs.\ Test Prep):

    \vspace{0.6em}
    Answer each with a \textbf{complete comparative sentence}:
    \begin{enumerate}\setlength{\itemsep}{0.55em}
      \item Compare the \textbf{spreads} (IQR) of the two groups.
            Name both groups; use a comparison word; give both IQRs; use context.
      \item The two IQR boxes overlap from 24 to 25. Is that a
            large or small overlap relative to the box widths?
            What does it suggest?
      \item Would a student scoring 32 be an outlier for the No Prep group?
            Show the fence calculation.
    \end{enumerate}

    \vspace{0.5em}
    {\color{navy}\small\textit{Key:
    IQRs are equal (both 6); small relative overlap suggests a real
    difference in center; No Prep upper fence $= 25 + 1.5(6) = 34 > 32$,
    not an outlier.}}
  \end{minipage}
\end{frame}
}

% ─────────────────────────────────────────────────────────────────────────────
% SLIDE 9 — GROUP ACTIVITY INTRO
% ─────────────────────────────────────────────────────────────────────────────
{
\setbeamercolor{background canvas}{bg=white}
\begin{frame}[t, plain]
  \begin{tikzpicture}[remember picture, overlay]
    \fill[navy] (current page.north west)
      rectangle ([yshift=-0.25cm]current page.north east);
  \end{tikzpicture}
  \vspace{0.45cm}\par
  \hspace{0.55cm}%
  \begin{minipage}{0.9\textwidth}
    \sectionlabel[greenacc]{Group Activity}
    \small
    \textbf{Hours of sleep for Athletes vs.\ Non-Athletes}\\[0.2em]
    \begin{tabular}{ll}
    Athletes ($n=15$):     & 6, 7, 7, 8, 8, 8, 9, 9, 9, 10, 10, 10, 11, 11, 12 \\
    Non-Athletes ($n=15$): & 5, 6, 6, 7, 7, 7, 7, 8, 8, 8, 9, 9, 10, 10, 14 \\
    \end{tabular}

    \vspace{0.5em}
    \textbf{Part 1} (8 min): Compute five-number summaries for each group.\\
    Find IQRs. Apply the fence rule: is 14 hours an outlier for Non-Athletes?\\
    What is the last non-outlier?

    \vspace{0.4em}
    \textbf{Part 2} (7 min): Read the parallel boxplots on your sheet.\\
    Verify your calculations. Answer interpretation questions.

    \vspace{0.4em}
    \textbf{Part 3} (5 min): Write a complete SOCS comparison.\\
    Name both groups in every sentence. Peer-review another group's response.
  \end{minipage}
\end{frame}
}

% ─────────────────────────────────────────────────────────────────────────────
% SLIDE 10 — EXIT TICKET
% ─────────────────────────────────────────────────────────────────────────────
{
\setbeamercolor{background canvas}{bg=sky}
\begin{frame}[t, plain]
  \navyheader{Exit Ticket}
  \vspace{0.8em}%
  \hspace{0.55cm}%
  \begin{minipage}{0.9\textwidth}
    \small
    \textit{Use the parallel boxplots on your exit ticket --- monthly rainfall
    for City A and City B.}

    \vspace{0.5em}
    \begin{enumerate}\setlength{\itemsep}{0.5em}
      \item Read the five-number summary for each city.
            Calculate both IQRs.
      \item Verify City B's outlier (10 in.) using the fence rule. Show work.
      \item Which city has a higher median? By how many inches?
      \item Compare the spreads (IQR). Which city is more consistent?
    \end{enumerate}

    \vspace{0.4em}
    \textbf{AP-Style:} Compare the distributions of monthly rainfall
    for City A and City B. Address shape, outlier(s), center, and spread.
    Use explicit comparative language naming both cities.
  \end{minipage}
\end{frame}
}

% ─────────────────────────────────────────────────────────────────────────────
% SLIDE 11 — LESSON SUMMARY
% ─────────────────────────────────────────────────────────────────────────────
{
\setbeamercolor{background canvas}{bg=navy}
\begin{frame}[plain]
  \begin{tikzpicture}[remember picture, overlay]
    \fill[navylight] (current page.north west)
      rectangle ([yshift=-1.35cm]current page.north east);
  \end{tikzpicture}
  \vspace{1.5cm}
  \hspace{0.55cm}%
  \begin{minipage}{0.9\textwidth}
    {\color{goldacc}\bfseries\large Lesson 1.9 Summary}\\[0.8cm]
    \small
    \begin{itemize}\setlength{\itemsep}{0.5em}
      \item {\color{white} \textbf{Parallel boxplots} on the same scale allow immediate
            comparison of center, spread, and shape across groups.}
      \item {\color{white} \textbf{AP grading} requires explicit comparative language:
            name both groups, use a comparison word, give both values.}
      \item {\color{white} A full \textbf{SOCS comparison} addresses shape, outliers,
            center (medians), and spread (IQRs) --- all in context.}
      \item {\color{white} A \textbf{back-to-back stemplot} shows individual values
            and shape for small datasets; read the median by counting to the
            middle leaf in each half.}
      \item {\color{skymid} \textbf{Next:} Lesson 1.10 --- The Normal distribution
            and $z$-scores: a mathematical model for symmetric, bell-shaped data.}
    \end{itemize}
  \end{minipage}
\end{frame}
}

\end{document}
